\chapter{Tensors, reduced-rank models, and inverse problems}\label{ch:inverse}
\begin{readergoal}Move from matrix factorization to multivariate measurements, and learn how identifiable models differ from underdetermined ones.\end{readergoal}
\section{More than two indices}
An experiment with neurons $n$, time $t$, and trial $r$ may be stored as a third-order tensor $\mathcal X_{ntr}$. Its modes represent different axes of variation. A CANDECOMP/PARAFAC (CP) approximation is
\[\mathcal X_{ntr}\approx\sum_{k=1}^K a_{nk}b_{tk}c_{rk}.\]
Interpretation requires confronting scaling and permutation ambiguities: multiplying one factor by $\alpha$ and dividing another by $\alpha$ leaves the product unchanged. Additional uniqueness conditions depend on the factors and rank. The mathematical pattern is broader than one physiological explanation.

\section{Reduced-rank regression}
For predictor matrix $X\in\mathbb R^{T\times p}$ and response matrix $Y\in\mathbb R^{T\times q}$, ordinary multivariate regression fits $Y\approx XB$. Reduced-rank regression requires $\rank(B)\le r$, controlling the dimension of predictive communication. Calling the resulting dimensions ``communication subspaces'' is an interpretation built on the predictive model, not proof that specific axons mediate the influence.

\section{The forward model is not the inverse model}
Let measurements obey $y=Hx+\varepsilon$, where $x$ is a source, $H$ a measurement operator, and $\varepsilon$ noise. The \emph{forward problem} computes $y$ from a chosen source and measurement physics. The \emph{inverse problem} estimates $x$ from $y$. If $H$ has a nontrivial kernel, $x$ and $x+z$ generate the same noiseless measurement whenever $Hz=0$. This is an identifiability failure, not a limitation remedied merely by better numerical precision.

For extracellular electric potentials, source geometry and volume conduction affect $H$, and several current patterns may explain the same electrode measurements. Biophysical treatments emphasize the interpretation of forward models and the assumptions behind inverse methods \cite{Buzsaki2012,Einevoll2013}. Regularization can choose one plausible solution, but it cannot magically restore missing independent information.

\section{Regularized reconstruction}
One possible estimator is $\hat x=\arg\min_x\{\|y-Hx\|_2^2+\lambda\|Lx\|_2^2\}$. If the normal equations yield an invertible matrix, $\hat x=(H^\top H+\lambda L^\top L)^{-1}H^\top y$. The operator $L$ expresses a prior preference (e.g., smoothness), while $\lambda$ sets its influence. The result answers a conditional question: which sources fit the observations well while obeying the chosen regularity assumption?

\begin{exercises}
\item Exhibit two distinct $x\in\mathbb R^2$ satisfying $[1\ 1]x=2$.
\item Explain why a CP decomposition's factors may be arbitrarily rescaled.
\item Describe what you would need to claim that a reduced-rank relationship reflects direct anatomical connectivity.
\end{exercises}